\documentclass[11pt]{beamer}

% Language and encoding settings
\usepackage[utf8]{inputenc}
\usepackage[T1]{fontenc}
\usepackage[english]{babel}

% Useful math and graphics packages
\usepackage{amsmath}
\usepackage{amsfonts}
\usepackage{amssymb}
\usepackage{graphicx}
\usepackage{booktabs} % For better looking tables

% Theme and color selection
\usetheme{Madrid} % Classic, professional theme
\usecolortheme{default}

% Informacje o prezentacji (strona tytułowa)
\title[]{Concept Drift in Media Bias Classification}
\subtitle{Machine Learning Methods for Data Streams}
\author[P. Jurczyk, K. Krawiec]{Piotr Jurczyk \and Krzysztof Krawiec}
\institute[WUT]
{
   \\
  Warsaw University of Technology
}
\begin{document}

% ---------------------------------------------------
% Title Slide
% ---------------------------------------------------
\begin{frame}
  \titlepage
\end{frame}

% ---------------------------------------------------
% Outline / Table of Contents
% ---------------------------------------------------
\begin{frame}{Outline}
  \tableofcontents
\end{frame}


\section{Introduction}

\begin{frame}{Introduction}
Modern online media can be treated as a \textbf{stream of textual data}, where political articles are published continuously and may reflect changing narratives, framing strategies, and ideological language patterns over time.

\vspace{0.3cm}

The aim of this project is to develop and evaluate two complementary approaches for \textbf{drift detection in political news streams}:
\begin{enumerate}
    \item \textbf{Part I}: Classification-based drift detection using online classifiers
    \item \textbf{Part II}: Unsupervised distributional drift detection
\end{enumerate}

\vspace{0.3cm}

Dataset: GDELT articles (Oct 2024 -- Apr 2026), 890 articles for binary left/right classification, 4200 articles in total.

\end{frame}

%---------------------------------------------------
\section{Synthetic Data Generation}

\begin{frame}{Synthetic Data Generation}

To validate drift detection under \textbf{controlled conditions}, we created a synthetic data generator with ideologically-coded vocabularies:

\vspace{0.3cm}

\begin{table}
\centering
\small
\begin{tabular}{lcc}
\toprule
\textbf{Component} & \textbf{Left} & \textbf{Right} \\
\midrule
Class vocabulary ($V_L$, $V_R$) & 50 terms & 50 terms \\
Shared neutral vocabulary & \multicolumn{2}{c}{130 terms} \\
Cross-contamination & \multicolumn{2}{c}{10\%} \\
Neutral noise & \multicolumn{2}{c}{40\%} \\
\bottomrule
\end{tabular}
\end{table}

\vspace{0.3cm}

\textbf{Drift simulation types:}
\begin{itemize}
    \item \textbf{Abrupt}: Instantaneous vocabulary swap
    \item \textbf{Gradual}: Distributed swap over time window
    \item \textbf{Recurring}: Vocabulary restoration
\end{itemize}

\end{frame}



%---------------------------------------------------
\section{Part I: Classification-Based Drift}

\begin{frame}{Part I: Classification-Based Drift}

\textbf{Approach:} Online classifiers process articles sequentially with test-then-train evaluation.

\vspace{0.3cm}

\textbf{Two vectorization strategies:}
\begin{enumerate}
    \item \textbf{TF-IDF (sparse)}: Lightweight bag-of-words baseline from \texttt{river}
    \item \textbf{LLM Embeddings (dense)}: all-MiniLM-L6-v2 Sentence Transformer (384-dim)
\end{enumerate}

\vspace{0.3cm}

\textbf{Classifiers:}
\begin{itemize}
    \item TF-IDF pipeline: MultinomialNB
    \item LLM pipeline: LogisticRegression (SGD)
    \item Baseline: Majority-class prediction (rolling 50-sample window)
\end{itemize}

\vspace{0.3cm}

\textbf{Drift detection:} ADWIN

\end{frame}

\begin{frame}{Synthetic Results: Basic Drift}

\begin{figure}
    \centering
    \includegraphics[width=1\textwidth]{Report/synthetic1.png}
\end{figure}

\vspace{-0.3cm}

\textbf{Key finding:} ADWIN drift detector successfully identifies all three injected drifts. Fresh models recover to near-100\% accuracy within hundreds of steps.

\end{frame}

\begin{frame}{Synthetic Results: Advanced Drift}

\begin{figure}
    \centering
    \includegraphics[width=1\textwidth]{Report/synthetic2.png}
\end{figure}

\vspace{-0.3cm}

\textbf{Key finding:} Abrupt and recurring drifts trigger ADWIN alerts. Gradual drift was not detected.

\end{frame}

\begin{frame}{Real-World Results: TF-IDF}

\begin{figure}
    \centering
    \includegraphics[width=1\textwidth]{Report/baseline.png}
\end{figure}

\vspace{-0.3cm}

\textbf{Critical finding:} Model accuracy overlaps perfectly with majority-class baseline.


\end{frame}

\begin{frame}{ Real-World Results: LLM Embeddings}

\begin{figure}
    \centering
    \includegraphics[width=1\textwidth]{Report/llm.png}
\end{figure}

\vspace{-0.3cm}

\textbf{Finding:} Dense embeddings also perform at baseline level with no improvement over majority-class prediction.


\end{frame}

%---------------------------------------------------
\section{Part II: Distributional Drift Detection}

\begin{frame}{Part II: Distributional Drift Detection}

\textbf{Approach:} Abandon requirement for accurate classification. Instead, detect drift by measuring semantic shifts directly.

\vspace{0.3cm}

\textbf{Pipeline:}
\begin{enumerate}
    \item \textbf{Aggregation}: Group articles into daily/weekly windows
    \item \textbf{Vectorization}: Embed aggregated documents
    \item \textbf{Distance stream}: Compute cosine distance between consecutive windows
    \item \textbf{Drift detection}: Apply ADWIN, PageHinkley, and KSWIN to distance stream
    \item \textbf{Event matching}: Validate detections against known political events (±7--14 days)
\end{enumerate}

\vspace{0.3cm}

\textbf{Advantages:}
\begin{itemize}
    \item No dependency on classification accuracy
    \item Directly measures semantic shifts in media discourse
    \item Multiple detectors improve robustness
\end{itemize}

\end{frame}

\begin{frame}{Part II Results: TF-IDF (Mixed Stream)}

\begin{figure}
    \centering
    \includegraphics[width=1\textwidth]{Report/h1_tfidf_daily_mixed.png}
\end{figure}

\vspace{-0.3cm}

\textbf{Observation:} TF-IDF produces clean drift detection with moderate event alignment around conventions and election period.

\end{frame}

\begin{frame}{Part II Results: Embeddings (Mixed Stream)}

\begin{figure}
    \centering
    \includegraphics[width=1\textwidth]{Report/h1_embedding_daily_mixed.png}
\end{figure}

\vspace{-0.3cm}

\textbf{Observation:} Dense embeddings exhibit noisier signals.

\end{frame}

\begin{frame}{Part II Results: Left-Leaning Stream}

\begin{figure}
    \centering
    \includegraphics[width=1\textwidth]{Report/h1_tfidf_daily_left.png}
\end{figure}

\vspace{-0.3cm}

\textbf{Key finding:} Sharp drift signals aligned with Democratic-specific events:
\begin{itemize}
    \item Biden's withdrawal (July 21, 2024)
    \item Harris nomination (August 5, 2024)
\end{itemize}

\end{frame}

\begin{frame}{Part II Results: Right-Leaning Stream}

\begin{figure}
    \centering
    \includegraphics[width=1\textwidth]{Report/h1_tfidf_daily_right.png}
\end{figure}

\vspace{-0.3cm}

\textbf{Key finding:} Pronounced drifts synchronized with Republican-specific events:
\begin{itemize}
    \item Trump conviction (May 30, 2024)
    \item Trump assassination attempt (July 13, 2024)
\end{itemize}

\end{frame}

%---------------------------------------------------
\section{Conclusions}

\begin{frame}{Summary of Findings}

\vspace{0.3cm}

\begin{block}{Part I: Classification-Based Drift Detection}
\textbf{Fails} on real-world data. Both TF-IDF and LLM approaches collapse to majority-class prediction, providing zero discriminative signal despite reliable performance on synthetic data.
\end{block}

\vspace{0.3cm}

\begin{block}{Part II: Distributional Drift Detection}
\textbf{Succeeds} by abandoning classification. TF-IDF outperforms embeddings. Ideology-specific stream analysis reveals clear event-drift alignment.
\end{block}


\end{frame}

\begin{frame}{Limitations \& Future Work}

\textbf{Limitations:}
\begin{itemize}
    \item Small corpus for classification
    \item Binary left/right categorization is oversimplified
    \item LLM embeddings not fine-tuned on political discourse
\end{itemize}

\vspace{0.3cm}

\textbf{Future directions:}
\begin{itemize}
    \item Larger, more diverse article corpus with multiple outlets per ideology
    \item Domain-specific fine-tuned embeddings (RoBERTa on political text)
\end{itemize}

\end{frame}

\begin{frame}
  \centering \Huge
  \emph{Thank you for your attention!}
  
  \vspace{1cm}
  \Large Any questions?
\end{frame}

\end{document}